% Thesis chapter: background and related work.
% Standalone chapter file; \input or \include from the thesis root.

\chapter{Background and Related Work}
\label{ch:background}

This chapter situates the thesis in four bodies of work: automation of
forestry cranes and comparable heavy machinery, robotic grasping in clutter,
learning for log grasping specifically, and the algorithmic building blocks
the thesis assembles (imitation learning, reinforcement learning fine-tuning,
and point cloud policy networks).

\section{Automation of Forestry Cranes and Heavy Machinery}

Mechanized timber harvesting has advanced steadily, but the machines remain
almost entirely manually operated, and the industry's own assessments identify
operator workload and scarcity as principal drivers for
autonomy~\cite{lindroos2017forestry, oliveira2021advances}. Classical
approaches automate the motion side of crane work. Semberg
et~al.~\cite{semberg2024} demonstrated automated single-log grasping on a
forwarder using one stereo camera for both detection and boom-tip visual
feedback control. Jebellat and Sharf~\cite{jebellat2023} generated crane
trajectories with dynamic programming to damp end-effector sway, and the
approach was subsequently implemented on a lab-scale manipulator emulating a
log loader~\cite{jebellat2025designing}. In adjacent heavy machinery,
reinforcement learning has been applied to hydraulic excavator
control~\cite{egli2022towards} and to material-handling machines with
underactuated, free-swinging tools~\cite{spinelli2024reinforcement}, and La
Hera and Morales~\cite{la2023autonomous} surveyed progress toward autonomous
forestry crane manipulation. These works automate motion generation for a
given target; they do not decide where to grasp in a dense pile, which is the
decision this thesis learns.

\section{Grasping in Clutter}

Learning-based grasp detection has progressed from isolated objects to
cluttered scenes. Deep grasp prediction networks map visual input directly to
grasp poses~\cite{morrison2018closing, kumra2017robotic}, and Dex-Net
demonstrated robust planning from synthetic point clouds at bin-picking
scale~\cite{mahler2017dexnet}. Sequential decluttering has been studied in bin
picking, including policies trained on simulated grasping sequences that
account for how each pick changes the heap~\cite{mahler2017sequential}, and
push-grasp synergies that actively create graspable
configurations~\cite{zeng2018learning}.

Two assumptions run through this literature and fail in the mill yard. First,
clutter is loose: objects are heterogeneous and separable, and the grasping
primitive is a precision or pinch grasp that isolates a single target, so the
core skill is singulation. Logs are near-uniform cylinders stacked in ordered,
tightly packed layers with extensive frictional contact; nothing can be
singulated, and the grapple captures bunches of up to twenty logs in one power
grasp. Second, the objective is usually the next grasp. In pile clearing the
objective is the whole sequence: each grasp causes settling, rolling, and
knock-offs that reshape the configuration for every subsequent attempt, and a
targeting strategy is only as good as its behavior over the full trajectory
from dense stack to empty rack.

\section{Learning for Log Grasping}

A growing body of work applies learning to log grasping, at smaller scales
than considered here. Wallin et~al.~\cite{wallin2024multilog} trained a
reinforcement learning agent with virtual visual servoing to grasp from small
groupings of two to five logs in simulation, reaching 95\% success. F{\"a}lldin
et~al.~\cite{falldin2025synthesizing} trained a U-Net on synthetic RGB-D
images to predict multi-log grasp poses for clusters of up to seven logs.
Ayoub et~al.~\cite{ayoub2023grasp} proposed a CNN grasp planner for a
log-loading forestry machine, validated in simulation and on a retrofitted
commercial loader. On the control side, Vu et~al.~\cite{vu2025autonomous}
released a MuJoCo forestry crane simulator and benchmarked RL velocity
controllers for grasping single logs, and Kurinov
et~al.~\cite{kurinov2025reinforcement} trained an RL agent for the full
single-log pick-and-place pipeline in simulation. Supporting perception work
includes log instance segmentation~\cite{fortin2022instance} and keypoint-based
grasp candidate detection~\cite{xu2021}.

Collectively these works address single-shot grasp prediction, single-log
manipulation, or small static groupings. None addresses sequential clearing of
a dense pile of hundreds of logs, where the pile state evolves under the
policy's own actions. To the author's knowledge, the work in this thesis is
the first study of learned sequential grasp targeting at this pile scale.

\section{Learning from Demonstrations and RL Fine-Tuning}

Behavior cloning, supervised learning of a policy from expert
state-action pairs, dates to ALVINN~\cite{pomerleau1988alvinn}. Its central
weakness is well understood: the student is graded on matching the expert's
action, not on the outcome of its own, so errors compound off the
demonstration distribution~\cite{ross2011dagger} and the student cannot exceed
the strategy it copies. Reinforcement learning optimizes outcomes directly but
explores poorly when reward is sparse and the observation is high-dimensional.
Combining the two, using demonstrations to initialize or regularize RL, has
proven effective in contact-rich manipulation~\cite{rajeswaran2018learning,
nair2018overcoming}. A complementary device for learning from raw sensing is
the asymmetric actor-critic~\cite{pinto2018asymmetric}: the critic, needed
only during training, consumes privileged simulator state and therefore learns
value quickly, while the actor trains on the deployable observation. This
thesis combines demonstration-initialized policies, on-policy fine-tuning with
proximal policy optimization~\cite{schulman2017ppo, schwarke2025rslrl}, and an
asymmetric critic, and contributes an analysis of what the combination
requires at pile scale: a frozen perception encoder, conservative trust
regions, and, ultimately, an action parameterization whose exploration cannot
leave the observed material.

\section{Policy Learning from Point Clouds}

Point clouds are unordered sets, so policy networks over them require
permutation-invariant architectures. PointNet~\cite{qi2017pointnet} applies a
shared per-point network followed by a symmetric max-pooling aggregation, and
its hierarchical successor PointNet++~\cite{qi2017pointnetplusplus} adds local
neighborhood structure. This thesis uses a PointNet encoder throughout, on
deliberately raw input: no segmentation or object detection is applied, and
the network must learn to attend to task-relevant geometry among logs, rack
structure, and ground. The choice is motivated in
Chapter~\ref{ch:sim_study}, where raw clouds match or exceed segmented ones
under behavior cloning, and stress-tested in Chapters~\ref{ch:sim2real} and
\ref{ch:bc_policies}, where the same property simplifies transfer to the real
sensor.
